\subsection{Who made the film: networks, stars, careers}
\begin{enumerate}[label={[\arabic*]},start=1,leftmargin=2.2em,itemsep=2pt]
\item Lash \& Zhao (2016). \textit{Early Predictions of Movie Success: The Who, What, and When of Profitability}. Journal of Management Information Systems, 33, 874--903. \url{https://arxiv.org/abs/1506.05382} \src{read (arXiv version)}
\item Packard et al. (2016). \textit{The role of network embeddedness in film success}. International Journal of Research in Marketing, 33, 328--342. \href{https://doi.org/10.1016/j.ijresmar.2015.06.007}{\nolinkurl{doi:10.1016/j.ijresmar.2015.06.007}} \src{read}
\item De Vany \& Walls (1999). \textit{Uncertainty in the movie industry: Does star power reduce the terror of the box office?} Journal of Cultural Economics, 23, 285--318. \href{https://doi.org/10.1023/A:1007608125988}{\nolinkurl{doi:10.1023/A:1007608125988}} \src{abstract}
\item Ravid (1999). \textit{Information, Blockbusters, and Stars: A Study of the Film Industry}. The Journal of Business, 72, 463--492. \url{https://ideas.repec.org/a/ucp/jnlbus/v72y1999i4p463-92.html} \src{abstract}
\item Elberse (2007). \textit{The Power of Stars: Do Star Actors Drive the Success of Movies?} Journal of Marketing, 71, 102--120.  \src{abstract}
\item Liu et al. (2018). \textit{Hot streaks in artistic, cultural, and scientific careers}. Nature, 559, 396--399. \href{https://doi.org/10.1038/s41586-018-0315-8}{\nolinkurl{doi:10.1038/s41586-018-0315-8}} \src{abstract}
\item John et al. (2017). \textit{Managerial ability and success: Evidence from the career paths of film directors}. Journal of Corporate Finance, 44, 425--439. \href{https://doi.org/10.1016/j.jcorpfin.2014.11.001}{\nolinkurl{doi:10.1016/j.jcorpfin.2014.11.001}} \src{abstract}
\item Cattani \& Ferriani (2008). \textit{A Core/Periphery Perspective on Individual Creative Performance: Social Networks and Cinematic Achievements in the {Hollywood} Film Industry}. Organization Science, 19, 824--844. \href{https://doi.org/10.1287/orsc.1070.0350}{\nolinkurl{doi:10.1287/orsc.1070.0350}} \src{read (first pages)}
\item Ferriani et al. (2009). \textit{The relational antecedents of project-entrepreneurship: Network centrality, team composition and project performance}. Research Policy, 38, 1545--1558. \href{https://doi.org/10.1016/j.respol.2009.09.001}{\nolinkurl{doi:10.1016/j.respol.2009.09.001}} \src{abstract}
\item Uzzi \& Spiro (2005). \textit{Collaboration and Creativity: The Small World Problem}. American Journal of Sociology, 111, 447--504. \href{https://doi.org/10.1086/432782}{\nolinkurl{doi:10.1086/432782}} \src{abstract}
\item Sharda \& Delen (2006). \textit{Predicting box-office success of motion pictures with neural networks}. Expert Systems with Applications, 30, 243--254.  \src{abstract}
\end{enumerate}
\subsection{What the film says: scripts, arcs, words}
\begin{enumerate}[label={[\arabic*]},start=12,leftmargin=2.2em,itemsep=2pt]
\item Del Vecchio et al. (2021). \textit{Improving productivity in {Hollywood} with data science: Using emotional arcs of movies to drive product and service innovation}. Journal of the Operational Research Society, 72, 1110--1137. \href{https://doi.org/10.1080/01605682.2019.1705194}{\nolinkurl{doi:10.1080/01605682.2019.1705194}} \src{read (arXiv 1807.02221 version)}
\item Reagan et al. (2016). \textit{The emotional arcs of stories are dominated by six basic shapes}. EPJ Data Science, 5, 31. \href{https://doi.org/10.1140/epjds/s13688-016-0093-1}{\nolinkurl{doi:10.1140/epjds/s13688-016-0093-1}} \src{abstract}
\item Toubia et al. (2021). \textit{How quantifying the shape of stories predicts their success}. Proceedings of the National Academy of Sciences, 118, e2011695118. \href{https://doi.org/10.1073/pnas.2011695118}{\nolinkurl{doi:10.1073/pnas.2011695118}} \src{read}
\item Berger et al. (2021). \textit{What Makes Content Engaging? How Emotional Dynamics Shape Success}. Journal of Consumer Research, 48, 235--250. \href{https://doi.org/10.1093/jcr/ucab010}{\nolinkurl{doi:10.1093/jcr/ucab010}} \src{read}
\item Eliashberg et al. (2007). \textit{From Story Line to Box Office: A New Approach for Green-Lighting Movie Scripts}. Management Science, 53, 881--893. \href{https://doi.org/10.1287/mnsc.1060.0668}{\nolinkurl{doi:10.1287/mnsc.1060.0668}} \src{read}
\item Eliashberg et al. (2014). \textit{Assessing Box Office Performance Using Movie Scripts: A Kernel-Based Approach}. IEEE Transactions on Knowledge and Data Engineering, 26, 2639--2648. \href{https://doi.org/10.1109/TKDE.2014.2306681}{\nolinkurl{doi:10.1109/TKDE.2014.2306681}} \src{abstract}
\item Hunter et al. (2016). \textit{Predicting box office from the screenplay: A text analytical approach}. Journal of Screenwriting, 7, 135--154. \href{https://doi.org/10.1386/josc.7.2.135_1}{\nolinkurl{doi:10.1386/josc.7.2.135_1}} \src{abstract}
\item Monroe et al. (2008). \textit{Fightin' Words: Lexical Feature Selection and Evaluation for Identifying the Content of Political Conflict}. Political Analysis, 16, 372--403. \href{https://doi.org/10.1093/pan/mpn018}{\nolinkurl{doi:10.1093/pan/mpn018}} \src{read}
\item Saxena \& Keller (2024). \textit{{MovieSum}: An Abstractive Summarization Dataset for Movie Screenplays}. Findings of the Association for Computational Linguistics: ACL 2024, 4043--4050. \url{https://aclanthology.org/2024.findings-acl.239} \src{abstract}
\end{enumerate}
\subsection{How it was received: critics, audiences, money}
\begin{enumerate}[label={[\arabic*]},start=21,leftmargin=2.2em,itemsep=2pt]
\item Holbrook \& Addis (2007). \textit{Taste versus the Market: An Extension of Research on the Consumption of Popular Culture}. Journal of Consumer Research, 34, 415--424.  \src{read}
\item Holbrook \& Addis (2008). \textit{Art versus commerce in the movie industry: A Two-Path Model of Motion-Picture Success}. Journal of Cultural Economics, 32, 87--107. \href{https://doi.org/10.1007/s10824-007-9059-2}{\nolinkurl{doi:10.1007/s10824-007-9059-2}} \src{abstract}
\item Holbrook (1999). \textit{Popular Appeal versus Expert Judgments of Motion Pictures}. Journal of Consumer Research, 26, 144--155. \href{https://doi.org/10.1086/209556}{\nolinkurl{doi:10.1086/209556}} \src{abstract}
\item Eliashberg \& Shugan (1997). \textit{Film Critics: Influencers or Predictors?} Journal of Marketing, 61, 68--78. \href{https://doi.org/10.2307/1251831}{\nolinkurl{doi:10.2307/1251831}} \src{abstract}
\item Basuroy et al. (2003). \textit{How Critical Are Critical Reviews? The Box Office Effects of Film Critics, Star Power, and Budgets}. Journal of Marketing, 67, 103--117. \href{https://doi.org/10.1509/jmkg.67.4.103.18692}{\nolinkurl{doi:10.1509/jmkg.67.4.103.18692}} \src{abstract}
\item Reinstein \& Snyder (2005). \textit{The Influence of Expert Reviews on Consumer Demand for Experience Goods: A Case Study of Movie Critics}. The Journal of Industrial Economics, 53, 27--51. \href{https://doi.org/10.1111/j.0022-1821.2005.00244.x}{\nolinkurl{doi:10.1111/j.0022-1821.2005.00244.x}} \src{abstract}
\item Moon et al. (2010). \textit{Dynamic Effects Among Movie Ratings, Movie Revenues, and Viewer Satisfaction}. Journal of Marketing, 74, 108--121.  \src{read (first pages)}
\item Basuroy et al. (2020). \textit{Is everybody an expert? An investigation into the impact of professional versus user reviews on movie revenues}. Journal of Cultural Economics, 44, 57--96.  \src{abstract}
\item Liu (2006). \textit{Word of Mouth for Movies: Its Dynamics and Impact on Box Office Revenue}. Journal of Marketing, 70, 74--89.  \src{abstract}
\item Duan et al. (2008). \textit{Do online reviews matter? An empirical investigation of panel data}. Decision Support Systems, 45, 1007--1016. \href{https://doi.org/10.1016/j.dss.2008.04.001}{\nolinkurl{doi:10.1016/j.dss.2008.04.001}} \src{abstract}
\item Joshi et al. (2010). \textit{Movie Reviews and Revenues: An Experiment in Text Regression}. Proceedings of NAACL-HLT 2010, 293--296. \url{https://aclanthology.org/N10-1038} \src{read}
\item Vig et al. (2012). \textit{The Tag Genome: Encoding Community Knowledge to Support Novel Interaction}. ACM Transactions on Interactive Intelligent Systems, 2, 13:1--13:44. \href{https://doi.org/10.1145/2362394.2362395}{\nolinkurl{doi:10.1145/2362394.2362395}} \src{abstract}
\end{enumerate}
\subsection{Methods and tools}
\begin{enumerate}[label={[\arabic*]},start=33,leftmargin=2.2em,itemsep=2pt]
\item Hutto \& Gilbert (2014). \textit{{VADER}: A Parsimonious Rule-Based Model for Sentiment Analysis of Social Media Text}. Proceedings of the International AAAI Conference on Web and Social Media (ICWSM), 216--225. 
\item Sanh et al. (2019). \textit{{DistilBERT}, a distilled version of {BERT}: smaller, faster, cheaper and lighter}. arXiv preprint arXiv:1910.01108. 
\item Blondel et al. (2008). \textit{Fast unfolding of communities in large networks}. Journal of Statistical Mechanics: Theory and Experiment, P10008. \href{https://doi.org/10.1088/1742-5468/2008/10/P10008}{\nolinkurl{doi:10.1088/1742-5468/2008/10/P10008}}
\item Kessler (2017). \textit{Scattertext: a Browser-Based Tool for Visualizing how Corpora Differ}. arXiv preprint arXiv:1703.00565. 
\end{enumerate}
